\documentclass[12pt]{article}
\usepackage[utf8]{inputenc}
\usepackage[T1]{fontenc}
\usepackage{hyperref}
\title{Zangskari Language Resources for K-12 Curriculum Development}
\author{Compiled by Claude Sonnet 5.5 (claude-sonnet-5-5) for Andr\'as Kornai}
\date{}
\begin{document}
\maketitle

\section{Introduction}
Zangskari (ISO 639-3 zau, Glottocode zang1248) is a Sino-Tibetan language of Zanskar in Ladakh, India, and is closely related to Ladakhi. Search results gave very little that is directly usable for schools. No published dictionary, grammar or school textbook turned up. This report lists only what was found.

\section{Primary Educational Resources}
No primers or textbooks were found in the search results.
A report on a seminar on documenting the Zanskari language mentions plans for a multilingual Zanskari--Hindi--Urdu--English dictionary, orthography development, primers and a pedagogical grammar. The report does not say that any of these exist yet.
\begin{itemize}
\item Seminar on documentation of Zanskari language held at GMDC Zanskar: \url{https://indusdispatch.in/seminar-on-documentation-of-zanskari-language-held-at-gmdc-zanskar/}
\end{itemize}

\section{Linguistic and Reference Materials}
\begin{itemize}
\item Wikipedia overview of the Zangskari language, covering dialect groups and classification: \url{https://en.wikipedia.org/wiki/Zangskari_language}
\item Endangered Languages Project sample: a narrative text with a glossary of words, definitions and pronunciations. This is a word list inside a story, not a dictionary: \url{https://endangeredlanguages.com/lang/3144/samples/9553}
\end{itemize}

\section{Religious and Community Resources}
No religious or community-organization resources specific to Zangskari were found in the search results.

\section{Digital and Audio Resources}
The CoRSAL collection in the UNT Digital Library holds materials documenting the Zangskari language and culture. This is the most promising source of audio and archival material.
\begin{itemize}
\item CoRSAL Zangskari collection (UNT Digital Library): \url{https://digital.library.unt.edu/explore/collections/ZSLR/}
\end{itemize}

\section{Curriculum Development Context}
Zangskari is endangered and is not known to be taught in schools. Any curriculum would have to be built from scratch. The CoRSAL archive and the planned documentation work are the likeliest starting points, and the people involved in that work should be contacted to confirm what exists. Speaker estimates vary widely between sources.

\end{document}